\documentclass[10pt,letterpaper,sans]{moderncv}
\moderncvstyle{banking}
\moderncvcolor{blue}

\usepackage[scale=0.81]{geometry}
\nopagenumbers{}
\setlength{\hintscolumnwidth}{2.25cm}

\name{Krishnaa}{Vadivel}
\title{Machine Learning Engineer | Scientific ML | Equivariant Models}
\email{srikrishnaa.careers@gmail.com}
\phone[mobile]{516-853-5663}
\social[linkedin]{sri-krishnaa-ece-phd}
\extrainfo{\href{https://github.com/krishnaa423}{GitHub: krishnaa423} \textbar{} \href{https://leetcode.com/u/krishnaa42342}{LeetCode: krishnaa42342} \textbar{} \href{https://hub.docker.com/u/krishnaa42342}{Docker Hub: krishnaa42342}}

\begin{document}
\makecvtitle

\section{Summary}
\cvitem{}{Machine-learning-focused PhD researcher working at the intersection of scientific computing, molecular/materials modeling, and scalable infrastructure. Experience spans PyTorch and equivariant graph models, agentic retrieval workflows, GPU acceleration, and distributed compute on national HPC systems.}

\section{Experience}
\cventry{2021--present}{PhD Research Assistant}{Yale University, Electrical Engineering}{}{}{\begin{itemize}%
\item Develop scientific ML workflows around first-principles materials research, including equivariant molecular models, physics-informed formulations, and ML-assisted acceleration of computational pipelines.
\item Build distributed and accelerator-enabled workflows in Python and C++ using CUDA, ROCm, MPI, OpenMP, and HPC environments such as Frontier and Perlmutter.
\item Explore theory-to-implementation links across graph neural networks, diffusion models, transformer-style abstractions, and data generation for scientific use cases.
\end{itemize}}
\cventry{2019--2021}{Software and Embedded Systems Engineer}{Probe Technologies Inc. / Weatherford Wireline Services}{}{}{\begin{itemize}%
\item Built GPU-accelerated processing and custom OpenGL visualization tools for large waveform datasets, providing practical experience with high-throughput data pipelines and performance-oriented engineering.
\item Supported production data and engineering workflows with ASP-style services and Microsoft SQL-backed tooling.
\end{itemize}}

\section{Selected Projects}
\cvitem{Equivariant Molecular ML}{PyTorch and graph-based workflows for molecular-property prediction and symmetry-aware learning on QM9-style data.}
\cvitem{Agentic Scientific RAG}{Retrieval-driven tooling that reads documentation and research papers to generate structured inputs and workflow scaffolds for computational science.}
\cvitem{Diffusion and Transformer Theory Notes}{Practical derivations and organized notes for DDPM-style models, conditioning, and index-notation ML formulations.}
\cvitem{Distributed Scientific Compute}{ML-adjacent acceleration work combining CUDA/ROCm, HPC scheduling, and scalable simulation pipelines.}

\section{Education}
\cventry{2021--present}{PhD, Electrical Engineering}{Yale University}{}{}{}
\cventry{2023}{MPhil, Electrical Engineering}{Yale University}{}{}{}
\cventry{2023}{MS, Electrical Engineering}{Yale University}{}{}{}
\cventry{2017--2018}{MEng, Electrical and Computer Engineering}{Cornell University}{}{}{}
\cventry{2013--2017}{BS, Electrical and Computer Engineering}{Cornell University}{}{}{}

\section{Technical Skills}
\cvitem{ML}{PyTorch, PyTorch Geometric, scikit-learn, equivariant graph neural networks}
\cvitem{Programming}{Python, C++, CUDA, JavaScript, Bash}
\cvitem{Data / Compute}{NumPy, pandas, SciPy, distributed workflows, CUDA, ROCm}
\cvitem{Systems}{Linux, Docker, Docker Compose, Kubernetes, gcloud}
\cvitem{Scientific}{PETSc, SLEPc, scientific simulation pipelines, physics-informed modeling}

\end{document}
